\hypertarget{target:chap3}{}\chapter{محاسبه مساحت با استفاده از انتگرال}\label{chap:area}
با استفاده از انتگرال معین می‌توان مساحت دقیق بین یک تابع و محور مختصات را در یک بازه مشخص به‌دست آورد.

\section{مساحت زیر نمودار توابع}
اگر تابع $f(x)$ در بازه $[a,b]$ پیوسته و مثبت باشد، مساحت ناحیه زیر نمودار آن از رابطه \eqref{eq:area_def} به‌دست می‌آید:
\begin{equation}
	A = \int_{a}^{b} f(x) \, dx.
	\label{eq:area_def}
\end{equation}

به‌عنوان مثال، اگر تابع $y=x^2$ را در بازه $[0,2]$ در نظر بگیریم، مساحت ناحیه با جایگذاری ضوابط در رابطه \eqref{eq:area_def}، به صورت فرآیند محاسباتی یکپارچه در رابطه \eqref{eq:area_calc} به‌دست می‌آید:
\begin{equation}
	A = \int_{0}^{2} x^2 \, dx = \left[ \frac{x^3}{3} \right]_0^2 = \frac{8}{3}.
	\label{eq:area_calc}
\end{equation}

همان‌طور که در رابطه \eqref{eq:area_calc} مشاهده می‌شود، حاصل این انتگرال‌گیری به یک مقدار عددی دقیق ختم شده است که مساحت زیر منحنی را نشان می‌دهد.

\section{الگوریتم و فلوچارت محاسباتی مجموع ریمان}
برای تقریب عددی مساحت زیر منحنی، الگوریتم \ref{alg:riemann_sum} فرآیند مجموع ریمان را پیاده‌سازی می‌کند. این الگوریتم در واقع ساختار پیوسته رابطه \eqref{eq:area_def} را به یک فرآیند گسسته تکرارپذیر تبدیل خواهد کرد.

\begin{latin}
	\begin{algorithm}
		\caption{Left Riemann Sum Approximation}
		\label{alg:riemann_sum}
		\begin{algorithmic}[1]
			\REQUIRE Function $f(x)$, interval $[a, b]$, number of subintervals $n$
			\ENSURE Approximated area $A$
			\STATE $\Delta x \gets \frac{b - a}{n}$
			\STATE $A \gets 0$
			\FOR{$i = 0$ \TO $n-1$}
			\STATE $x \gets a + i \cdot \Delta x$
			\STATE $A \gets A + f(x) \cdot \Delta x$
			\ENDFOR
			\RETURN $A$
		\end{algorithmic}
	\end{algorithm}
\end{latin}

فلوچارت این فرآیند با استفاده از ابزار تیکز\footnote{\lr{TikZ}} در شکل \ref{fig:flowchart_area} به تصویر کشیده شده است.

\begin{figure}[htbp]
	\centering
	\begin{tikzpicture}[node distance=1.5cm, auto]
		\tikzset{
			startstop/.style={draw, rectangle, rounded corners, minimum width=3cm, minimum height=0.8cm, text centered, draw=blue!70, fill=blue!10, thick},
			io/.style={draw, trapezium, trapezium left angle=70, trapezium right angle=110, minimum width=3cm, minimum height=0.8cm, text centered, draw=orange!70, fill=orange!10, thick},
			process/.style={draw, rectangle, minimum width=3.5cm, minimum height=0.8cm, text centered, draw=green!70, fill=green!10, thick},
			decision/.style={draw, diamond, aspect=1.8, minimum width=3cm, minimum height=0.8cm, text centered, draw=red!70, fill=red!10, thick, inner sep=0pt},
			arrow/.style={thick,->,>=stealth}
		}
		
		\node (start) [startstop] {شروع};
		\node (input) [io, below=0.6cm of start] {ورودی: $a, b, n, f(x)$};
		\node (init) [process, below=0.6cm of input] {$\Delta x = \frac{b-a}{n}$ , $A=0$ , $i=0$};
		\node (cond) [decision, below=1cm of init] {$i < n$؟};
		\node (calc) [process, left=1.5cm of cond] {$A = A + f(a + i\Delta x)\Delta x$};
		\node (incr) [process, above=1cm of calc] {$i = i + 1$};
		\node (stop) [startstop, below=1cm of cond] {خروجی $A$ و پایان};
		
		\draw [arrow] (start) -- (input);
		\draw [arrow] (input) -- (init);
		\draw [arrow] (init) -- (cond);
		\draw [arrow] (cond) -- node[anchor=south] {بله} (calc);
		\draw [arrow] (calc) -- (incr);
		
		% تغییر اصلی اینجاست: اتصال مستقیم و مورب بین دو گره
		\draw [arrow] (incr.south east) -- (cond.north west);
		
		\draw [arrow] (cond) -- node[anchor=east] {خیر} (stop);
	\end{tikzpicture}
	\caption{فلوچارت مراحل اجرای الگوریتم مجموع ریمان برای مساحت}
	\label{fig:flowchart_area}
\end{figure}

در جدول \ref{Table1}، نمادهای پرکاربرد ریاضی استفاده شده در محاسبات با ویژگی‌های خواسته شده معرفی گردیده‌اند.

\begin{table}[ht]
	\centering
	\caption{نمادهای استفاده شده در محاسبات انتگرالی}
	\label{Table1}
	\renewcommand{\arraystretch}{1.5}
	\setlength{\tabcolsep}{10pt}
	\rowcolors{2}{blue!10}{white}
	\begin{tabular}{c!{\color{blue}\vrule width 1.5pt}c}
		\hline 
		نماد & توضیح \\  
		\noalign{\hrule height 1.5pt}
		$f(x)$ & تابع مورد نظر \\
		$a, b$ & کران‌های انتگرال \\
		$A$ & مساحت \\
		$V$ & حجم \\
		$dx$ & عنصر دیفرانسیلی \\
		\hline
	\end{tabular}
\end{table}